\documentclass[10pt]{article}

% Specify the margins and text size.
\setlength{\textwidth}{6.8in}
\setlength{\textheight}{9.5in}
\setlength{\oddsidemargin}{0pt}
\setlength{\evensidemargin}{0pt}
\setlength{\topmargin}{0pt}
\setlength{\hoffset}{-.1in}
\setlength{\voffset}{-.8in}

\setlength{\parskip}{5pt}
\setlength{\parindent}{0pt}

% Load some fonts and symbol packages
\usepackage{latexsym}
\usepackage{pifont}       % contains 'star' symbol for counterinsurgency handbook title
\usepackage{yfonts} 
\usepackage{amsmath}
\usepackage{amsfonts}

\usepackage{graphicx}     % actually, this is loaded by pstricks
\usepackage[T1]{fontenc}
\usepackage{ifthen}
\usepackage{pstricks,pst-func,pst-text,pst-node,multido,pst-plot,calc,pst-3dplot}
%\usepackage[all]{xy}
%\usepackage{animate}

% The hyperref package inserts links into the pdf.
\definecolor{MyLinkColor}{rgb}{.1,.2,1}
\definecolor{MyCiteColor}{rgb}{.1,1,.2}
\definecolor{MyURLColor}{rgb}{.4,.4,.4}
\usepackage[backref=true,pagebackref=false,hyperindex,colorlinks=true,
  linkcolor=MyLinkColor,urlcolor=MyURLColor]{hyperref}

% The tweaklist package is something I found on the web.  It provides a simple interface
% for making changes to spacing used in the itemize and enumerate environments.  Comment
% this out if you don't care to use tweaklist.
\usepackage{tweaklist}
\renewcommand{\itemhook}{\setlength{\parskip}{2pt}\setlength{\parsep}%
{1pt}\setlength{\topsep}{0pt}\setlength{\itemsep}{0pt}}

\newcommand{\N}{\ensuremath\mathbb{N}}
\newcommand{\Q}{\ensuremath\mathbb{Q}}
\newcommand{\R}{\ensuremath\mathbb{R}}
\newcommand{\C}{\ensuremath\mathbb{C}}
\newcommand{\z}{\hphantom{0}}

\begin{document}
\pagestyle{empty}

\href{http://www.su.edu}{\includegraphics[height=1.75cm]{sulogo.eps}}
\vspace{-1.69cm}

{\small{\ }\hfill
\begin{tabular}{cl}
& Math 207\\& Introduction to Statistics\\
\end{tabular}
}
\setlength{\baselineskip}{1.05\baselineskip}
\medskip

{\setlength{\baselineskip}{1.05\baselineskip}
\begin{center}
\textbf{\large  Worksheet:  Introduction to Descriptive Statistics}
\end{center}


1. Summer camp attendees are asked to rate food quality at the camp.  Here are the results:\vspace{-3pt}
\begin{center}
Good, Good, Good, Fair, Good, Fair, Poor, Good, Poor, Fair, Good, Good, Fair, Good, Fair, Fair\vspace{-4pt}
\end{center}
Make a frequency table then a bar plot of the data.  
\vspace{2.5in}

2. Over many years, final class averages are compiled for a math class.  The results are tabulated below.

\begin{tabular}{|cc|}\hline
Score  & Percent\vphantom{\LARGE V}\\\hline
\z0 -- 60  & \z5\vphantom{\LARGE Y}\\
60 -- 70 & 10\\
70 -- 80 & 23\\
80 -- 90 & 35\\
\z90 -- 100 & 27\\\hline
\end{tabular}
%\end{center}

\hspace{10pt} a) Draw a histogram of the data.
\vspace{.5in}

\hspace{10pt} b) If a student scored exactly 90\%, in which bin would the score be classified?  (trick question)
\vspace{.4in}

\hspace{10pt} c) About what percent of students scored between 90 and 93?
\vspace{.4in}

3. A computer simulation of flipping coins is being tested.  The results are tabulated below.\vspace{-4pt}
\begin{center}
\begin{tabular}{|cc|}\hline
Number of Coin Flips & Percent of Heads\vphantom{\LARGE V}\\\hline
\hphantom{,000}10                  & 80\vphantom{\LARGE V}\\
\hphantom{,00}100                  & 55\\
\hphantom{,0}1000                  & 57\\
10,000                             & 57\\\hline
\end{tabular}\vspace{-4pt}
\end{center}
Is the simulation working correctly?  Which row of the table gives the strongest evidence to support your claim?

\vfill
\eject
\end{document}
